% =====================================================================
%  CASSA Observatory -- Photometric Pipeline: Participant Handbook
%
%  Build (run twice, so the table of contents settles):
%      pdflatex participant_handbook.tex
%      pdflatex participant_handbook.tex
%
%  Needs a TeX install with sourceserifpro, sourcesanspro, inconsolata and
%  tcolorbox -- all in texlive-latex-recommended / texlive-fonts-extra.
% =====================================================================
\documentclass[11pt,a4paper]{article}

\usepackage[T1]{fontenc}
\usepackage[utf8]{inputenc}
\usepackage{textcomp}
\usepackage[default,semibold]{sourceserifpro}
\usepackage[semibold]{sourcesanspro}
\usepackage[varqu,varl]{inconsolata}
\usepackage{microtype}

\usepackage[a4paper,top=2.3cm,bottom=2.3cm,left=2.5cm,right=2.5cm,headsep=14pt]{geometry}
\usepackage{xcolor}
\usepackage{titlesec}
\usepackage{enumitem}
\usepackage{booktabs}
\usepackage{tabularx}
\usepackage{array}
\usepackage{longtable}
\usepackage{fancyhdr}
\usepackage{fancyvrb}
\usepackage[breakable,skins]{tcolorbox}
\usepackage{parskip}
\usepackage[hidelinks]{hyperref}

% ---------------------------------------------------------------- palette
\definecolor{ink}{HTML}{111820}
\definecolor{inksoft}{HTML}{54606E}
\definecolor{inkfaint}{HTML}{7C8996}
\definecolor{accent}{HTML}{1C6B8B}
\definecolor{accentsoft}{HTML}{E6F1F6}
\definecolor{rule}{HTML}{D2DCE4}
\definecolor{sunk}{HTML}{EEF2F6}
\definecolor{pass}{HTML}{2C7A55}
\definecolor{warn}{HTML}{9E6614}
\definecolor{fail}{HTML}{AF382A}
\definecolor{warnsoft}{HTML}{FBF2E2}

\color{ink}
\hypersetup{colorlinks=true,linkcolor=accent,urlcolor=accent,
            pdftitle={CASSA Photometric Pipeline -- Participant Handbook},
            pdfauthor={CASSA Observatory}}

% ---------------------------------------------------------------- headings
\titleformat{\section}
  {\sffamily\bfseries\Large\color{ink}}
  {\color{accent}\normalsize\thesection\hspace{0.7em}}{0pt}{}
  [\vspace{2pt}{\color{ink}\titlerule[1.2pt]}]
\titlespacing*{\section}{0pt}{22pt}{10pt}

\titleformat{\subsection}{\sffamily\bfseries\normalsize\color{ink}}{}{0pt}{}
\titlespacing*{\subsection}{0pt}{14pt}{4pt}

\titleformat{\subsubsection}[runin]
  {\sffamily\bfseries\small\color{inksoft}}{}{0pt}{}[\hspace{0.5em}]
\titlespacing*{\subsubsection}{0pt}{8pt}{6pt}

\setlist[itemize]{leftmargin=1.1em,itemsep=2pt,topsep=3pt,parsep=0pt}
\setlist[enumerate]{leftmargin=1.4em,itemsep=2pt,topsep=3pt,parsep=0pt}

% ---------------------------------------------------------------- running head
\pagestyle{fancy}
\fancyhf{}
\renewcommand{\headrulewidth}{0.4pt}
\renewcommand{\headrule}{\hbox to\headwidth{\color{rule}\leaders\hrule height \headrulewidth\hfill}}
\fancyhead[L]{\sffamily\scriptsize\color{inkfaint}CASSA Photometric Pipeline \textbullet\ Participant Handbook}
\fancyhead[R]{\sffamily\scriptsize\color{inkfaint}\thepage}

% ---------------------------------------------------------------- code
\DefineVerbatimEnvironment{cmdtext}{Verbatim}
  {fontsize=\small,fontfamily=tt,commandchars=\\\{\},xleftmargin=0pt}

\newtcolorbox{cmd}{
  breakable, enhanced, sharp corners, boxrule=0pt, leftrule=2pt,
  colback=sunk, colframe=accent, boxsep=0pt,
  left=9pt, right=9pt, top=7pt, bottom=7pt,
  before skip=8pt, after skip=10pt
}

\newcommand{\cm}[1]{\texttt{\small #1}}
\newcommand{\cmt}[1]{\textcolor{inkfaint}{#1}}

% ---------------------------------------------------------------- callouts
\newtcolorbox{note}[1][Note]{
  breakable, enhanced, sharp corners, boxrule=0pt, leftrule=2.5pt,
  colback=accentsoft, colframe=accent,
  left=10pt, right=10pt, top=8pt, bottom=8pt,
  before skip=10pt, after skip=12pt,
  fonttitle=\sffamily\bfseries\scriptsize,
  coltitle=inksoft, title={\MakeUppercase{#1}}, attach title to upper=\par\vspace{3pt}
}
\newtcolorbox{caution}[1][Careful]{
  breakable, enhanced, sharp corners, boxrule=0pt, leftrule=2.5pt,
  colback=warnsoft, colframe=warn,
  left=10pt, right=10pt, top=8pt, bottom=8pt,
  before skip=10pt, after skip=12pt,
  fonttitle=\sffamily\bfseries\scriptsize,
  coltitle=warn, title={\MakeUppercase{#1}}, attach title to upper=\par\vspace{3pt}
}

% notebook entry
\newtcolorbox{nbbox}[2]{
  breakable, enhanced, sharp corners=all, boxrule=0.5pt,
  colback=white, colframe=rule,
  left=12pt, right=12pt, top=10pt, bottom=10pt,
  before skip=10pt, after skip=12pt,
  title={\sffamily\bfseries\small\colorbox{accent}{\textcolor{white}{\ \texttt{#1}\ }}\hspace{6pt}#2},
  coltitle=ink, colbacktitle=white, titlerule=0pt,
  attach boxed title to top left={xshift=0pt,yshift*=0pt},
  boxed title style={empty}
}

\newcommand{\verdict}[3]{%
  \noindent\colorbox{#1!14}{\textcolor{#1}{\sffamily\bfseries\scriptsize\ #2\ }}%
  \hspace{6pt}\begin{minipage}[t]{0.80\textwidth}\small #3\end{minipage}\par\vspace{5pt}}

\newcommand{\recordline}[2]{%
  \vspace{6pt}{\color{rule}\hrule}\vspace{6pt}%
  \noindent{\sffamily\bfseries\scriptsize\color{accent}\MakeUppercase{#1}}\hspace{6pt}%
  \begin{minipage}[t]{0.78\textwidth}\small #2\end{minipage}}

\setlength{\tabcolsep}{8pt}
\renewcommand{\arraystretch}{1.25}
\newcommand{\thead}[1]{{\sffamily\bfseries\scriptsize\color{inkfaint}\MakeUppercase{#1}}}

\begin{document}

% =====================================================================
%  Masthead
% =====================================================================
\thispagestyle{empty}
\begin{flushleft}
{\sffamily\bfseries\scriptsize\color{accent}\MakeUppercase{CASSA Observatory \textbullet\ Hands-on session}}\par
\vspace{10pt}
{\sffamily\bfseries\fontsize{28}{32}\selectfont Photometric Pipeline\\Participant Handbook}\par
\vspace{14pt}
{\large\color{inksoft} Everything you need for the session: install the pipeline on your own
machine, put a night of observatory data where it belongs, and work through five notebooks
that take raw frames to a calibrated magnitude --- with an error bar you can defend.}\par
\vspace{20pt}
{\color{rule}\hrule height 1pt}
\vspace{10pt}
\begin{tabularx}{\textwidth}{@{}XXXX@{}}
\thead{You will produce} & \thead{Notebooks} & \thead{Disk needed} & \thead{Python} \\
{\sffamily\small one star's mag $\pm$ error} & {\sffamily\small 5, run in order} &
{\sffamily\small $\sim$1\,GB} & {\sffamily\small 3.10 or newer} \\
\end{tabularx}
\vspace{6pt}
{\color{rule}\hrule height 1pt}
\end{flushleft}

\vspace{18pt}
{\sffamily\bfseries\large Contents}\vspace{-14pt}
\begin{small}
\renewcommand{\contentsname}{}
\setcounter{tocdepth}{1}
\tableofcontents
\end{small}

\clearpage

% =====================================================================
\section{Before you start}
% =====================================================================
Do the install \textbf{before} the session, not during it. Building the environment downloads a
few hundred megabytes and takes ten to twenty minutes on a normal connection; doing it in the
room costs everyone the first hour.

\vspace{4pt}
\begin{tabularx}{\textwidth}{@{}p{6.2cm}X@{}}
\toprule
\thead{What} & \thead{Why} \\
\midrule
A laptop running Linux or macOS & One install route for both, on Intel and ARM alike. Windows
works too, but through WSL --- so start that part early, it needs a reboot. \\
$\sim$1\,GB free disk & 267\,MB for the night's raw frames, $\sim$325\,MB for what the
pipeline writes, plus the environment. On Windows, allow another $\sim$2\,GB for WSL. \\
Internet, at least once & Installing, and Phase 3's reference-catalog query. Both cache, so a
second pass runs offline. \\
Comfort with a terminal & You will type about six commands. No prior astronomy software needed. \\
\bottomrule
\end{tabularx}

\begin{note}[Bring this]
Your assigned star --- the instructor gives you either a catalog \cm{NUMBER} or an RA/Dec pair.
Measuring it is the session's deliverable.
\end{note}

% =====================================================================
\section{Install}
% =====================================================================
One route, on every operating system: get \textbf{Miniforge}, then run one
script. Windows gets there through WSL --- see below, and do that part first.

\begin{note}[Why Miniforge and not plain pip]
Phase 2 plate-solves your images against a star catalog, and every solver that
can do it is a compiled binary rather than a Python package. Miniforge is what
supplies one, along with prebuilt versions of the whole scientific stack, so
nothing has to compile on your laptop.
\end{note}

\subsection{Which platform you are on}

\begin{tabularx}{\textwidth}{@{}p{4.6cm}p{3.6cm}X@{}}
\toprule
\thead{Your machine} & \thead{Route} & \thead{Notes} \\
\midrule
Linux, x86-64            & \cm{./install.sh}        & The reference platform. \\
Linux, ARM (aarch64)     & \cm{./install.sh -{}-conda} & Works, but only with ASTAP --- nothing else is published for ARM Linux. \\
macOS, Intel             & \cm{./install.sh}        & Identical to Linux. \\
macOS, Apple Silicon     & \cm{./install.sh}        & Identical. See the Gatekeeper note below. \\
Windows                  & \cm{.\textbackslash install.ps1} & Works natively now. WSL also fine if you already use it. \\
\bottomrule
\end{tabularx}

\subsection{Windows}
Windows installs natively. Get Miniforge (or use a Miniconda you already have),
then from PowerShell in the cloned repository:

\begin{cmd}
\begin{cmdtext}
.\textbackslash{}install.ps1
cassa-doctor
\end{cmdtext}
\end{cmd}

It does what \cm{install.sh} does on Linux and macOS: builds the environment,
installs the pipeline, fetches ASTAP --- the one plate solver published for
Windows --- registers the Jupyter kernel, and checks its own work.

\begin{note}[WSL still works, if you prefer it]
WSL is real x86-64 Linux, so every instruction in this handbook applies
unchanged inside it. \cm{.\textbackslash install.ps1 -Wsl} prints the setup
steps. If you go that way, keep the repository and the work directory in your
WSL home --- your Windows drives appear under \cm{/mnt/c}, and reducing a night
across \cm{/mnt} is several times slower.
\end{note}

\subsection{Step 1 --- Install Miniforge}
Miniforge is a minimal conda that defaults to the conda-forge channel. If you
already have conda or mamba, skip this step.

The installer name is built from your own machine, so this is the same pair of
commands on Linux, on macOS (Intel and Apple Silicon alike), and inside WSL:

\begin{cmd}
\begin{cmdtext}
curl -L -O "https://github.com/conda-forge/miniforge/releases/latest/\textbackslash{}
download/Miniforge3-$(uname)-$(uname -m).sh"
bash Miniforge3-$(uname)-$(uname -m).sh
\end{cmdtext}
\end{cmd}

Accept the defaults, then \textbf{close and reopen your terminal} so \cm{conda}
is on your PATH. Check it:

\begin{cmd}
\begin{cmdtext}
conda --version
\end{cmdtext}
\end{cmd}

On macOS you also need Apple's command-line tools, once, for \cm{git}:

\begin{cmd}
\begin{cmdtext}
xcode-select --install
\end{cmdtext}
\end{cmd}

\subsection{Step 2 --- Clone and install}
Identical everywhere:

\begin{cmd}
\begin{cmdtext}
git clone https://github.com/cassaiub/observatory.git
cd observatory/photometric_pipeline
./install.sh
\end{cmdtext}
\end{cmd}

The script builds a conda environment called \cm{cassa-photometry} from
\cm{environment.yml}, installs the pipeline into it, adds whichever plate solver
your platform can run, and finishes by checking its own work. Expect ten to
twenty minutes and a few hundred megabytes.

On ARM Linux add \cm{-{}-conda}: the in-process solver publishes no ARM wheel,
so the conda route plus ASTAP is the combination that works there.

\subsection{Step 3 --- Activate and verify}

\begin{cmd}
\begin{cmdtext}
conda activate cassa-photometry
cassa-doctor
\end{cmdtext}
\end{cmd}

\cm{cassa-doctor} is the one command that proves the install: it reports your
platform, the package version, the environment, which plate solver you have and
which one a run would use, and where the sky data will come from.
\textbf{Run it before the session, and bring the output if anything is not
green.} You need to run \cm{conda activate cassa-photometry} in every new
terminal.

\begin{note}[Check the kernel name before you start]
The installer registers a Jupyter kernel called \textbf{Python (CASSA
photometry)}. When a notebook opens, check that name appears in the top-right
corner; if it does not, use \emph{Kernel $\rightarrow$ Change Kernel}.

This is the single most common way the session goes wrong, and it looks like a
broken install rather than a wrong kernel: the notebook reports
\cm{ModuleNotFoundError: No module named 'cassa\_photometry'} on a machine where
the install clearly succeeded. Installing packages sets up an \emph{environment};
it does not by itself make Jupyter use it.
\end{note}

\begin{note}[Which solver you get, and why it does not matter]
The installer tries three, in order, and the first that installs wins:
\textbf{ASTAP} (a sub-megabyte binary that works on every platform, including
ARM), then Astrometry.net's \cm{solve-field}, then an in-process Python solver.
\cm{cassa-doctor} tells you which one you have.

\textbf{Your results do not depend on the answer.} The one thing ASTAP does not
report --- the astrometric residual \cm{ASTRMS}, which Phase 3 uses to size its
catalog cross-match --- the pipeline measures for itself, so every backend
produces it. If your neighbour has a different solver, your catalogs are still
comparable.
\end{note}

\begin{note}[There is nothing to download in advance]
Plate solving matches your stars against a sky catalog, and those catalogs are
large --- 859\,MB for ASTAP, some 34\,GB for Astrometry.net. You need neither.
The pipeline reads your frame's pointing and field size, works out which pieces
that field needs, and fetches only those: about \textbf{6\,MB} for ASTAP or
165\,MB for Astrometry.net, cached under \cm{\textasciitilde/.cache/cassa-photometry/}.
Pointing somewhere new later costs only the pieces the new field adds.
\end{note}

\begin{caution}[macOS: Gatekeeper may quarantine ASTAP]
If \cm{cassa-doctor} finds the ASTAP binary but a solve fails with a security
warning, macOS has quarantined the downloaded file. Clear it once:
\cm{xattr -d com.apple.quarantine "\$(command -v astap\_cli)"}
\end{caution}
% =====================================================================
\section{Where the data goes}
% =====================================================================
The acquisition software writes a night as a tree, sorted by frame type and then by target.
Copy that whole dated folder, unchanged, into \cm{photometric\_pipeline/workshop/raw/}:

\begin{cmd}
\begin{cmdtext}
workshop/raw/
`-- 20260903                    \cmt{<- the night, as the software named it}
    |-- BIAS/untargeted/          \cmt{15 frames}
    |-- DARK/untargeted/          \cmt{15 frames}
    |-- FLAT/untargeted/          \cmt{56 frames}
    `-- LIGHT/m22/                \cmt{20 frames}
        `-- 20260903T043707.171323_m22_Luminance_LIGHT_iub-dhaka_....fits
\end{cmdtext}
\end{cmd}

Do not rename, flatten or sort anything. \cm{RAW\_DIR} is searched \textbf{recursively}, so the
tree is read exactly as delivered --- and a single flat directory of frames works too, if that
is what you are given.

\begin{note}[Folders do not decide anything]
Every frame is classified by its \cm{IMAGETYP}, \cm{FILTER} and \cm{EXPTIME} headers, never by
the folder it sits in. A frame filed under the wrong folder is still reduced as whatever it
actually is --- notebook 01 tells you when the two disagree.
\end{note}

\subsection{Filters in this night}
\textcolor[HTML]{3E6FB8}{\sffamily\bfseries\small Blue}\quad
\textcolor[HTML]{3F8A55}{\sffamily\bfseries\small Green}\quad
\textcolor[HTML]{B4483A}{\sffamily\bfseries\small Red}\quad
\textcolor[HTML]{6C7885}{\sffamily\bfseries\small Luminance}\par
\vspace{4pt}
Five science frames in each, at 60\,s. Flats exist for all four.

\subsection{What the pipeline writes}
Everything lands under \cm{workshop/work/}, one directory per phase. Re-running a phase
replaces its products in place rather than piling up copies.

\vspace{4pt}
\begin{tabularx}{\textwidth}{@{}p{2.9cm}p{2.5cm}X@{}}
\toprule
\thead{Directory} & \thead{Written by} & \thead{Contents} \\
\midrule
\cm{work/phase1} & notebook 02 & \cm{calibrated\_*.fits} --- SCI/ERR/DQ, in electrons \\
\cm{work/phase2} & notebook 03 & \cm{Master\_*.fits} --- stacked and WCS-solved \\
\cm{work/phase3} & notebook 04 & \cm{*\_catalog.csv}, \cm{*\_fluxcal.fits}, segmentation maps \\
\bottomrule
\end{tabularx}

\subsection{If your data lives somewhere else}
Point at it instead of copying --- the notebooks honour these environment variables:

\begin{cmd}
\begin{cmdtext}
export RAW_DIR=/path/to/the/night
export WORK_DIR=/path/for/outputs
\end{cmdtext}
\end{cmd}

% =====================================================================
\section{Running the notebooks}
% =====================================================================
Activate the environment, then start Jupyter from the notebooks directory:

\begin{cmd}
\begin{cmdtext}
conda activate cassa-photometry
cd observatory/photometric_pipeline/workshop/notebooks
jupyter lab
\end{cmdtext}
\end{cmd}

\textbf{Run them in order.} Notebook 00 checks your environment and 01 reads only the raw
frames, but 02, 03 and 04 each consume the previous notebook's output. Within a notebook, run
every cell from the top --- later cells depend on names defined earlier.

\subsection{The first cell of every notebook}
Each notebook opens with the same config cell. It resolves \cm{RAW\_DIR}, \cm{WORK\_DIR} and
the per-phase directories, then prints what it found --- including your raw tree, directory by
directory. If those paths look wrong, stop and fix them there; everything downstream inherits
them.

\subsection{Filling in the blanks}
The exercises are partly written for you. Lines you must complete are marked:

\begin{cmd}
\begin{cmdtext}
gain = FILL_IN        \cmt{# TODO 1: e-/ADU as Phase 1 resolved it}
\end{cmdtext}
\end{cmd}

The surrounding code is scaffolding --- leave it alone. Each blank takes one expression. If you
run a cell without filling them you get \cm{NameError: name 'FILL\_IN' is not defined}, which
names the exact line still waiting for you. Nothing fails silently.

\subsection{Reading the answer tables}
Every exercise heading carries a small table naming the one or two values it wants, with the
expected answer beside each. Compare your numbers against it --- a mismatch is worth raising in
the room, because it usually means something real about your data rather than a mistake in your
arithmetic.

% =====================================================================
\section{Notebook by notebook}
% =====================================================================

\begin{nbbox}{00}{Setup \quad\textnormal{\sffamily\scriptsize\color{inkfaint}$\sim$5 min}}
\small\color{inksoft} Four checks that your machine can do the session at all. Run every cell.
\par\vspace{6pt}\color{ink}
{\sffamily\bfseries\small What you are checking}
\begin{itemize}
\item The package imports and reports a version.
\item The astrometry index directory resolves and holds index files.
\item \cm{solve-field} is found --- or is not, in which case the in-process solver takes over,
which is fine.
\item Your raw tree is found, and the frame counts per directory match what you copied in.
\end{itemize}
\recordline{Stop if}{the raw tree shows zero frames, or the index directory does not exist.
Neither is recoverable later.}
\end{nbbox}

\begin{nbbox}{01}{Raw frame health \quad\textnormal{\sffamily\scriptsize\color{inkfaint}$\sim$10 min}}
\small\color{inksoft} Are these frames worth reducing? The pipeline will happily reduce bad
calibration frames and produce plausible numbers that are quietly meaningless --- this is where
that gets caught.
\par\vspace{6pt}\color{ink}
{\sffamily\bfseries\small What it does}
\begin{itemize}
\item Inventories every frame by type, filter and exposure, and flags any frame whose folder
disagrees with its header.
\item Measures each frame once --- level, scatter, saturated fraction, temperature --- and plots
every frame's signal level against the ADC ceiling.
\item Runs sixteen checks and prints a verdict.
\end{itemize}
\vspace{6pt}
{\sffamily\bfseries\small The verdict}\par\vspace{5pt}
\verdict{pass}{PROCEED}{Everything passed. Go on to 02.}
\verdict{warn}{CAUTION}{Nothing blocking, but the result will be weaker than it looks. Note the
warnings --- they explain later numbers.}
\verdict{fail}{STOP}{A blocking problem. Each one names a way the reduction will be
\emph{wrong} rather than merely noisy, and none can be repaired afterwards. Raise it with an
instructor.}
\recordline{Record}{the verdict, and the headline of any FAIL.}
\end{nbbox}

\begin{nbbox}{02}{Phase 1 --- Calibration \quad\textnormal{\sffamily\scriptsize\color{inkfaint}$\sim$20 min}}
\small\color{inksoft} Bias, dark and flat come off; the frame is converted to electrons; and the
error budget is seeded. Output is \texttt{calibrated\_*.fits} carrying SCI, ERR and DQ.
\par\vspace{6pt}\color{ink}
The run itself takes several minutes for a full night --- start it, then read the exercise text
while it works.
\vspace{6pt}

{\sffamily\bfseries\small Exercises}
\begin{itemize}
\item[] {\ttfamily\scriptsize\color{inkfaint}Ex 1 \textbullet\ 3 blanks}\quad
\textbf{What did calibration actually change?} Find the instrumental pedestal that was removed,
and the real sky left behind. Both in electrons.
\item[] {\ttfamily\scriptsize\color{inkfaint}Ex 2 \textbullet\ 2 blanks}\quad
\textbf{Is the ERR plane just Poisson $+$ read noise?} Find the ratio of actual to predicted
error, and what the excess is.
\end{itemize}
\recordline{Record}{four numbers: pedestal, sky, ratio, extra term.}
\end{nbbox}

\begin{nbbox}{03}{Phase 2 --- Integration \& WCS \quad\textnormal{\sffamily\scriptsize\color{inkfaint}$\sim$20 min}}
\small\color{inksoft} Frames are aligned, stacked with inverse-variance weighting, and
plate-solved against real stars. Output is \texttt{Master\_*.fits} with a WCS.
\par\vspace{6pt}\color{ink}
{\sffamily\bfseries\small Exercises}
\begin{itemize}
\item[] {\ttfamily\scriptsize\color{inkfaint}Ex 1 \textbullet\ 2 blanks}\quad
\textbf{How much deeper is the stack?} Find the depth gain and the number of frames stacked.
Compare against $\sqrt{N}$.
\item[] {\ttfamily\scriptsize\color{inkfaint}Ex 2 \textbullet\ 2 blanks}\quad
\textbf{What did the plate solve buy us?} Find the plate scale in arcsec/pixel and the field of
view in arcmin.
\end{itemize}
\recordline{Watch for}{a depth gain that falls short of $\sqrt{N}$. That is real, not an
error: warping resamples pixels, which correlates neighbours, and correlated noise does not
average down as fast as the formula assumes.}
\end{nbbox}

\begin{nbbox}{04}{Phase 3 --- Photometry \& zero point \quad\textnormal{\sffamily\scriptsize\color{inkfaint}$\sim$25 min}}
\small\color{inksoft} Sources are detected and measured, a zero point is derived by
cross-matching against a reference catalog, and a catalog with magnitudes and errors is
written. This is the notebook that needs internet.
\par\vspace{6pt}\color{ink}
{\sffamily\bfseries\small Exercises}
\begin{itemize}
\item[] {\ttfamily\scriptsize\color{inkfaint}Ex 1 \textbullet\ 2 blanks}\quad
\textbf{Check the zero-point arithmetic.} Recompute magnitude and error from flux and the
header; find the largest residual, and which term dominates for the brightest star.
\item[] {\ttfamily\scriptsize\color{inkfaint}Ex 2 \textbullet\ 2 blanks}\quad
\textbf{Your assigned star.} Fill in \texttt{ASSIGNED\_NUMBER} or \texttt{ASSIGNED\_RADEC} at
the top of the cell, then report the magnitude, its uncertainty, and which half of the error
budget dominates.
\end{itemize}
\end{nbbox}

\begin{caution}[If there is no zero point]
Phase 3 writes \cm{\_fluxcal.fits} only when it measures one. Without it the catalog still
exists, carrying instrumental magnitudes, and the notebook falls back to those automatically
--- every shape you are asked about survives, only the zero of the scale is arbitrary.
Exercise 1 skips itself and says so.
\end{caution}

% =====================================================================
\section{What to hand in}
% =====================================================================
One line, from Exercise 2 of notebook 04:

\vspace{4pt}
\begin{tabularx}{\textwidth}{@{}p{5.4cm}X@{}}
\toprule
\thead{Item} & \thead{Example} \\
\midrule
Which star & \cm{NUMBER 7, RA 278.98621, Dec -23.79178} \\
Magnitude $\pm$ uncertainty & \cm{13.756 $\pm$ 0.017} \\
Dominant error term & \cm{zero point} \\
Filter and frames stacked & \cm{Blue, 3 frames} \\
\bottomrule
\end{tabularx}

\begin{note}[Quote it properly]
A magnitude without an uncertainty is not a measurement, and one written to more digits than
its uncertainty supports is a claim you cannot defend. Round the value to the precision the
error allows --- the notebook does this for you, and it is worth understanding why.
\end{note}

% =====================================================================
\section{When something breaks}
% =====================================================================
\begin{longtable}{@{}p{5.6cm}p{9.2cm}@{}}
\toprule
\thead{What you see} & \thead{What it means \& what to do} \\
\midrule
\endhead
\cm{NameError: name 'FILL\_IN'\newline is not defined} &
A blank you have not filled yet. The traceback names the line. \\
\cm{NameError} on a name from the config cell &
Your kernel imported the config before it was last edited. \textbf{Kernel $\rightarrow$
Restart}, then run from the first cell. \\
Notebook 01 says \textbf{STOP} &
A calibration frame is unusable. Do not reduce past it --- the numbers will look fine and be
wrong. Raise it in the room. \\
\cm{solver: solve-field}\newline \cm{not on PATH} &
Not a fault. Three backends can solve, and \cm{cassa-doctor}'s \cm{solver: in use} line says
which one you have. Only a problem if \cm{solver: any} \emph{fails} --- that means none is
installed. \\
\cm{no star database covering}\newline \cm{this field} &
ASTAP could not fetch the sky tiles for your pointing. Check your connection; the tiles are
about 6\,MB and are fetched during the solve. \\
macOS: ASTAP found, but the\newline solve fails with a security\newline warning &
Gatekeeper has quarantined the downloaded binary. Clear it once with
\cm{xattr -d com.apple.quarantine "\$(command -v astap\_cli)"} \\
Phase 3: \emph{zero point failed}, no \cm{\_fluxcal.fits} &
The cross-match found no usable stars in your image. The catalog is still written with
instrumental magnitudes and the exercises adapt. Worth investigating, not worth blocking on. \\
Phase 3 hangs or times out &
It is querying an online reference catalog. Check your connection; the query caches, so a
second run is fast. \\
\cm{conda not found} &
Miniforge is not installed, or you did not reopen the terminal after installing it. On Windows,
make sure you are inside the Ubuntu (WSL) window, not PowerShell. \\
A phase takes far longer than the estimate &
Normal on a laptop for a full night. Watch the progress bars; they move per frame. On WSL,
also check the repository is in your Linux home rather than under \cm{/mnt/c} --- that alone
can cost several times the runtime. \\
\bottomrule
\end{longtable}

Still stuck? \cm{cassa-doctor} prints the state of the whole install in one screen. Bring that
output rather than a description.

\vspace{18pt}
{\color{rule}\hrule}
\vspace{8pt}
{\sffamily\scriptsize\color{inkfaint}
CASSA Observatory \textbullet\ Independent University, Bangladesh --- rooftop observatory,
Dhaka.\\
Pipeline documentation lives in the repository: \texttt{README.md},
\texttt{docs/REFERENCE.md} and \texttt{docs/CUSTOMIZING.md}.}

\end{document}
